% Options for packages loaded elsewhere
% Options for packages loaded elsewhere
\PassOptionsToPackage{unicode}{hyperref}
\PassOptionsToPackage{hyphens}{url}
\PassOptionsToPackage{dvipsnames,svgnames,x11names}{xcolor}
%
\documentclass[
  11pt,
  a4paper]{article}
\usepackage{xcolor}
\usepackage[left=1in,top=1in,right=1in,bottom=1in]{geometry}
\usepackage{amsmath,amssymb}
\setcounter{secnumdepth}{-\maxdimen} % remove section numbering
\usepackage{iftex}
\ifPDFTeX
  \usepackage[T1]{fontenc}
  \usepackage[utf8]{inputenc}
  \usepackage{textcomp} % provide euro and other symbols
\else % if luatex or xetex
  \usepackage{unicode-math} % this also loads fontspec
  \defaultfontfeatures{Scale=MatchLowercase}
  \defaultfontfeatures[\rmfamily]{Ligatures=TeX,Scale=1}
\fi
\usepackage{lmodern}
\ifPDFTeX\else
  % xetex/luatex font selection
\fi
% Use upquote if available, for straight quotes in verbatim environments
\IfFileExists{upquote.sty}{\usepackage{upquote}}{}
\IfFileExists{microtype.sty}{% use microtype if available
  \usepackage[]{microtype}
  \UseMicrotypeSet[protrusion]{basicmath} % disable protrusion for tt fonts
}{}
\makeatletter
\@ifundefined{KOMAClassName}{% if non-KOMA class
  \IfFileExists{parskip.sty}{%
    \usepackage{parskip}
  }{% else
    \setlength{\parindent}{0pt}
    \setlength{\parskip}{6pt plus 2pt minus 1pt}}
}{% if KOMA class
  \KOMAoptions{parskip=half}}
\makeatother
% Make \paragraph and \subparagraph free-standing
\makeatletter
\ifx\paragraph\undefined\else
  \let\oldparagraph\paragraph
  \renewcommand{\paragraph}{
    \@ifstar
      \xxxParagraphStar
      \xxxParagraphNoStar
  }
  \newcommand{\xxxParagraphStar}[1]{\oldparagraph*{#1}\mbox{}}
  \newcommand{\xxxParagraphNoStar}[1]{\oldparagraph{#1}\mbox{}}
\fi
\ifx\subparagraph\undefined\else
  \let\oldsubparagraph\subparagraph
  \renewcommand{\subparagraph}{
    \@ifstar
      \xxxSubParagraphStar
      \xxxSubParagraphNoStar
  }
  \newcommand{\xxxSubParagraphStar}[1]{\oldsubparagraph*{#1}\mbox{}}
  \newcommand{\xxxSubParagraphNoStar}[1]{\oldsubparagraph{#1}\mbox{}}
\fi
\makeatother


\usepackage{longtable,booktabs,array}
\newcounter{none} % for unnumbered tables
\usepackage{calc} % for calculating minipage widths
% Correct order of tables after \paragraph or \subparagraph
\usepackage{etoolbox}
\makeatletter
\patchcmd\longtable{\par}{\if@noskipsec\mbox{}\fi\par}{}{}
\makeatother
% Allow footnotes in longtable head/foot
\IfFileExists{footnotehyper.sty}{\usepackage{footnotehyper}}{\usepackage{footnote}}
\makesavenoteenv{longtable}
\usepackage{graphicx}
\makeatletter
\newsavebox\pandoc@box
\newcommand*\pandocbounded[1]{% scales image to fit in text height/width
  \sbox\pandoc@box{#1}%
  \Gscale@div\@tempa{\textheight}{\dimexpr\ht\pandoc@box+\dp\pandoc@box\relax}%
  \Gscale@div\@tempb{\linewidth}{\wd\pandoc@box}%
  \ifdim\@tempb\p@<\@tempa\p@\let\@tempa\@tempb\fi% select the smaller of both
  \ifdim\@tempa\p@<\p@\scalebox{\@tempa}{\usebox\pandoc@box}%
  \else\usebox{\pandoc@box}%
  \fi%
}
% Set default figure placement to htbp
\def\fps@figure{htbp}
\makeatother


% definitions for citeproc citations
\NewDocumentCommand\citeproctext{}{}
\NewDocumentCommand\citeproc{mm}{%
  \begingroup\def\citeproctext{#2}\cite{#1}\endgroup}
\makeatletter
 % allow citations to break across lines
 \let\@cite@ofmt\@firstofone
 % avoid brackets around text for \cite:
 \def\@biblabel#1{}
 \def\@cite#1#2{{#1\if@tempswa , #2\fi}}
\makeatother
\newlength{\cslhangindent}
\setlength{\cslhangindent}{1.5em}
\newlength{\csllabelwidth}
\setlength{\csllabelwidth}{3em}
\newenvironment{CSLReferences}[2] % #1 hanging-indent, #2 entry-spacing
 {\begin{list}{}{%
  \setlength{\itemindent}{0pt}
  \setlength{\leftmargin}{0pt}
  \setlength{\parsep}{0pt}
  % turn on hanging indent if param 1 is 1
  \ifodd #1
   \setlength{\leftmargin}{\cslhangindent}
   \setlength{\itemindent}{-1\cslhangindent}
  \fi
  % set entry spacing
  \setlength{\itemsep}{#2\baselineskip}}}
 {\end{list}}
\usepackage{calc}
\newcommand{\CSLBlock}[1]{\hfill\break\parbox[t]{\linewidth}{\strut\ignorespaces#1\strut}}
\newcommand{\CSLLeftMargin}[1]{\parbox[t]{\csllabelwidth}{\strut#1\strut}}
\newcommand{\CSLRightInline}[1]{\parbox[t]{\linewidth - \csllabelwidth}{\strut#1\strut}}
\newcommand{\CSLIndent}[1]{\hspace{\cslhangindent}#1}



\setlength{\emergencystretch}{3em} % prevent overfull lines

\providecommand{\tightlist}{%
  \setlength{\itemsep}{0pt}\setlength{\parskip}{0pt}}



 


\usepackage{iftex}

% % Minion Pro
% \usepackage[lf]{MinionPro}
% \usepackage{MyriadPro}

% IBM Plex
% \usepackage{fontspec}
% \usepackage{unicode-math}
% \setmainfont{IBM Plex Serif}
% \setsansfont{IBM Plex Sans}
% \setmonofont{IBM Plex Mono}
% \setmathfont{IBM Plex Math Beta 240212}

% STIX Two
% \usepackage{fontspec}
% \setmainfont{Stix Two Text}
% \setmathfont{Stix Two Math}
% \setsansfont[Scale=0.9608140457]{Fira Sans} % 7.2489/7.54454

% Libertinus Serif + Libertinus Math + Source Sans Pro + Source Code Pro
\ifPDFTeX
\usepackage[mono=false,osf]{libertine}
\usepackage[libertine]{newtxmath}
% \usepackage[scale=1,osf]{sourcesanspro} % 7.1832/7.1832
% \usepackage[scale=0.8772671637]{sourcecodepro} % 7.1832/7.36934*0.9
\usepackage[scale=0.9181263229,osf]{roboto} % 7.1832/7.82376
\usepackage[scale=0.8132241635]{roboto-mono} % 7.1832/7.94969*0.9
\else
\usepackage[mono=false,osf]{libertine}
\setmathfont[Scale=MatchUppercase]{libertinusmath-regular.otf}
% \usepackage[scale=0.9850312868,osf]{sourcesanspro} % 7.2051/7.31459
% \usepackage[scale=0.8865281581]{sourcecodepro} % 7.2051/7.31459*0.9
\usepackage[scale=0.9126225152,osf]{roboto} % 7.2051/7.89494
\usepackage[scale=0.8213602637]{roboto-mono} % 7.2051/7.89494*0.9
\fi

% Enable section and paragraph numbering
\setcounter{secnumdepth}{4} % Numbering down to paragraph level
\setcounter{tocdepth}{4} % Include paragraphs in the ToC

%%% CUSTOMIZE SECTION SIZES
\makeatletter
\renewcommand{\section}{%
  \@startsection{section}{1}
    {0pt} 					% indent
    {1.0em plus 0.2em minus 0.2em} 	% beforeskip
    {0.001em plus 0.001em minus 0.001em}    	% afterskip
    {\Large\bfseries\sffamily}%
}
\renewcommand{\subsection}{%
  \@startsection{subsection}{2}
    {0pt} 					% indent
    {1.0em plus 0.2em minus 0.2em} 	% beforeskip
    {0.001em plus 0.001em minus 0.001em}    	% afterskip
    {\large\bfseries\sffamily}%
}
\renewcommand{\subsubsection}{%
  \@startsection{subsubsection}{3}
    {0pt} 					% indent
    {1.0em plus 0.2em minus 0.2em} 	% beforeskip
    {0.001em plus 0.001em minus 0.001em}    	% afterskip
    {\normalsize\bfseries\sffamily}%
}
\renewcommand{\paragraph}{%
  \@startsection{paragraph}{4}
    {0pt} 					% indent
    {1.0em plus 0.2em minus 0.2em} 	% beforeskip
    {0.001em plus 0.001em minus 0.001em}    	% afterskip
    {\normalsize\bfseries\sffamily}%
}
\makeatother

%%% GENERAL PACKAGES
\usepackage{hyphenat} % don't hyphenat titles
\usepackage{graphicx}
\usepackage{tabularx} 
\usepackage{marvosym} % Mail logo
\usepackage{enumitem} % Enumerate items (list)
\usepackage{amsmath} % for enhanced math features
\usepackage[labelfont=bf,justification=justified,singlelinecheck=false]{caption} % Figures caption setup
\usepackage[rightcaption]{sidecap} % figure with sidecaption
\usepackage[]{lineno} % Line numbers
\usepackage{ragged2e} % justify
\usepackage{array}
\usepackage{lscape} % landscape tables
\usepackage{rotating} % landscape tables
\usepackage{bm} % bold math

%%%% ORCID LOGO
\usepackage{academicons}
\definecolor{orcidlogocol}{HTML}{A6CE39}
\usepackage{orcidlink}

%%%% LAYOUT 
\tolerance=200
\emergencystretch=2em
\hyphenpenalty=1000
\hbadness=2000
\widowpenalty=1000 % widow penalty
\clubpenalty=1000 % orphans penalty
\renewcommand{\baselinestretch}{1.3333333333333333333333333333} % line spacing

\renewcommand{\arraystretch}{1.3333}
\makeatletter
\@ifpackageloaded{float}{}{\usepackage{float}}
\floatstyle{plain}
\@ifundefined{c@chapter}{\newfloat{appfig}{h}{loappfig}}{\newfloat{appfig}{h}{loappfig}[chapter]}
\floatname{appfig}{Figure A}
\newcommand*\quartoappfigref[1]{Figure \hyperref[#1]{A\ref{#1}}}
\@ifpackageloaded{caption}{}{\usepackage{caption}}
\DeclareCaptionLabelFormat{quartoappfigreflabelformat}{#1#2}
\captionsetup[appfig]{labelformat=quartoappfigreflabelformat}
\newcommand*\listofappfigs{\listof{appfig}{List of Appendix Figures}}
\makeatother
\makeatletter
\@ifpackageloaded{caption}{}{\usepackage{caption}}
\AtBeginDocument{%
\ifdefined\contentsname
  \renewcommand*\contentsname{Table of contents}
\else
  \newcommand\contentsname{Table of contents}
\fi
\ifdefined\listfigurename
  \renewcommand*\listfigurename{List of Figures}
\else
  \newcommand\listfigurename{List of Figures}
\fi
\ifdefined\listtablename
  \renewcommand*\listtablename{List of Tables}
\else
  \newcommand\listtablename{List of Tables}
\fi
\ifdefined\figurename
  \renewcommand*\figurename{Figure}
\else
  \newcommand\figurename{Figure}
\fi
\ifdefined\tablename
  \renewcommand*\tablename{Table}
\else
  \newcommand\tablename{Table}
\fi
}
\@ifpackageloaded{float}{}{\usepackage{float}}
\floatstyle{ruled}
\@ifundefined{c@chapter}{\newfloat{codelisting}{h}{lop}}{\newfloat{codelisting}{h}{lop}[chapter]}
\floatname{codelisting}{Listing}
\newcommand*\listoflistings{\listof{codelisting}{List of Listings}}
\makeatother
\makeatletter
\makeatother
\makeatletter
\@ifpackageloaded{caption}{}{\usepackage{caption}}
\@ifpackageloaded{subcaption}{}{\usepackage{subcaption}}
\makeatother
\usepackage{bookmark}
\IfFileExists{xurl.sty}{\usepackage{xurl}}{} % add URL line breaks if available
\urlstyle{same}
\hypersetup{
  pdftitle={What is the influence of knee joint movement on the maximal average mechanical power output of human quadriceps femoris muscle?},
  pdfauthor={Edwin D.H.M. Reuvers; A.J. `Knoek' van Soest; Dinant A. Kistemaker},
  colorlinks=true,
  linkcolor={blue},
  filecolor={Maroon},
  citecolor={Blue},
  urlcolor={Blue},
  pdfcreator={LaTeX via pandoc}}


\title{What is the influence of knee joint movement on the maximal
average mechanical power output of human quadriceps femoris muscle?}
\author{Edwin D.H.M. Reuvers \and A.J. `Knoek' van Soest \and Dinant A.
Kistemaker}
\date{2026-04-13}
\begin{document}

\begin{titlepage}
% This is a combination of Pandoc templating and LaTeX
% Pandoc templating https://pandoc.org/MANUAL.html#templates
% See the README for help
\begin{titlepage}

\raggedright
% Title and subtitle
{\sffamily\LARGE\bfseries\nohyphens{What is the influence of knee joint
movement on the maximal average mechanical power output of human
quadriceps femoris muscle?}}\\[2\baselineskip]


% Authors
% This hairy bit of code is just to get "and" between the last 2 authors.
	    \mbox{\normalsize{Edwin D.H.M.
Reuvers}\textsuperscript{1}\textsuperscript{,}{\textsuperscript{ \Letter}} \orcidlink{https://orcid.org/0000-0001-8932-1959},}
  % Below is nice format,  but then it gives spaces where I do not want it!
	%  	\mbox{{\normalsize{Edwin D.H.M. Reuvers}}
	% 		% 		% 	  \textsuperscript{1}
	% 		% 		% 	  	% 	    \textsuperscript{,}
	% 	  {\textsuperscript{ \Letter}}
	% 		% 		% 	  \orcidlink{https://orcid.org/0000-0001-8932-1959}
	% 	}
	    \mbox{\normalsize{A.J. `Knoek' van Soest}\textsuperscript{1},}
  % Below is nice format,  but then it gives spaces where I do not want it!
	%  	\mbox{{\normalsize{A.J. `Knoek' van Soest}}
	% 		% 		% 	  \textsuperscript{1}
	% 		% 		% 	}
				\mbox{and \normalsize{Dinant A.
Kistemaker}\textsuperscript{1} \orcidlink{https://orcid.org/0000-0003-1535-009X}}
		% Below is nice format,  but then it gives spaces where I do not want it!
		% \mbox{
		%  and \normalsize{Dinant A. Kistemaker}
		% 		% 		%   \textsuperscript{1}
		% 		% 		% 		%   \orcidlink{https://orcid.org/0000-0003-1535-009X}, 
		% 		% }
	
\vspace{0.5\baselineskip}

%%%%%% Affiliations & Corresponding author
  \begin{enumerate}[label={~\arabic*}, leftmargin=*, align=left, labelsep=0em, nosep]
  %
  
                  \item[\small\textsuperscript{1}]
        \small{Faculty of Behavioural and Movement Sciences, Vrije
Universiteit Amsterdam, Amsterdam Movement Sciences, The Netherlands,}
        \small{Department of Human Movement Sciences}
            
  \vspace{0.5\baselineskip}
  
            \item[\small\textsuperscript{\Letter}]
      \small{Corresponding author:~Edwin D.H.M.
Reuvers (\href{mailto:edwin.reuvers.academia@pm.me}{edwin.reuvers.academia@pm.me})}
                    \end{enumerate}

\vspace{1\baselineskip}

%%%%%% Project website
  \small{Access the code and full analysis notebook at \url{https://edwinreuvers.github.io/publications/hq-ampo}}

\vspace{1\baselineskip}

%%%%%% Abstract
  \begin{abstract}
    \justifying
    Performance during motor activities such as wheelchair riding,
    cycling, rowing and speed skating, depends critically on the average
    mechanical power output (AMPO) produced by the muscles. To maximise
    short-duration performance, limb movements should allow muscles to
    deliver maximal AMPO. However, it is unclear which movement
    maximises AMPO of human muscle. In this study, we employed a
    Hill-type muscle-tendon-complex (MTC) model to predict the maximally
    attainable AMPO of human \emph{m. quadriceps femoris} for various
    imposed periodic knee joint movements. Based on these predictions,
    we selected one set of conditions predicted to yield identical
    maximally attainable AMPO despite substantial variations in knee
    joint movements and another set of conditions predicted to yield
    substantial variations in maximally attainable AMPO. In the
    experiment, periodic knee joint movements were fully imposed by a
    knee dynamometer. Participants were instructed to maximise AMPO and,
    to this end, received visual feedback on their cumulative mechanical
    work throughout each cycle. Experimental data closely matched
    predictions derived from the Hill-type MTC model, confirming the
    validity of the model. Model predictions showed a substantial
    influence of knee joint movement on the maximally attainable AMPO.
    Specifically, predictions revealed a strong interaction between
    cycle frequency and knee joint excursion: increasing one
    necessitates a decrease in the other to maximise AMPO. Even more
    interestingly, \emph{m. quadriceps femoris} should spend about 80\%
    of the cycle duration while shortening, independent of cycle
    frequency and/or knee joint excursion.
  \end{abstract}

\vfill
\vspace{0.1\textheight}
\end{titlepage}\end{titlepage}

\linenumbers

\section{Introduction}\label{introduction}

Empty

Grieve et al. (\citeproc{ref-grieve_prediction_1978}{1978})

NOT HTML

\section*{Acknowledgments}\label{acknowledgments}
\addcontentsline{toc}{section}{Acknowledgments}

Tester

\section{Appendix}\label{appendix}

\subsection{Appendix 1}\label{appendix-1}

\subsection{Appendix 2}\label{appendix-2}

\section{Appendix}\label{appendix-3}

\subsection{Appendix 1}\label{appendix-1-1}

\section*{References}\label{references}
\addcontentsline{toc}{section}{References}

\protect\phantomsection\label{refs}
\begin{CSLReferences}{1}{1}
\bibitem[\citeproctext]{ref-grieve_prediction_1978}
\textbf{Grieve, D. W., Pheasant, S. and Cavanagh, P. R.} (1978).
Prediction of gastrocnemius length from knee and ankle joint posture. In
\emph{Biomechanics {VI}-{A}} (ed. Asmussen, E.) and Jorgensen, K.), pp.
405--412. Baltimore: University Park Press.

\end{CSLReferences}




\end{document}
